\documentclass{article} %[a4paper,10pt]
\usepackage[utf8]{inputenc}

\usepackage[english]{babel}

\usepackage{indentfirst}
\usepackage{multicol}
\usepackage{bbm}
\usepackage{amsmath}
\usepackage{amssymb}
% \usepackage{mathrsfs}
\usepackage{physics}
% \usepackage{calrsfs}
\usepackage{mathtools}
\usepackage{mathpunctspace}
\usepackage{wrapfig}

\usepackage[makeroom]{cancel}

\usepackage{graphicx}  % To include figures
\usepackage{subcaption}  % To use subfigure or subcaption environments


\usepackage{booktabs}

\usepackage{tikz}
\usepackage{pgfplots}
\usepackage{pgfplotstable}

\pgfplotsset{
    every axis plot/.append style={line width=0.5pt}
}

\usepackage{float}
\usepackage{hyperref}
\newfloat{graph}{htbp}{grp}
\floatname{graph}{Graf}

% Bibtex
\usepackage[
    %backend=biber, 
    natbib=true,
    style=numeric,
    sorting=none
]{biblatex}
\addbibresource{bibliography.bib} %Import the bibliography file
\usepackage{csquotes}

% \usepackage{algpseudocode}
\usepackage{algorithm, algorithmic}

\usepackage{xcolor}
\usepackage{graphicx}
\graphicspath{{graphics/}{../output/}}
\usepackage[top = 1cm, bottom = 1cm, left = 1cm, right = 1cm]{geometry}

\setlength{\textfloatsep}{1pt plus 1.0pt minus 2.0pt}

\title{Allen-Cahn model}
\author{Ján Kovačovský}


\pgfplotsset{compat=1.18}


\newcommand{\todo}[1]{\textcolor{red}{#1}}
\newcommand{\note}[1]{\textcolor{blue}{#1}}

\setcounter{section}{3}


\begin{document}

\section*{Allen-Cahn model}

The Allen-Cahn model describes the motion of a diffuse interface as the $L^2$ gradient flow of the Ginzburg-Landau energy \cite{AllenCahn1979}
\begin{equation}
    \mathcal{E}(\phi) = \int_{\Omega} \frac{1}{4 \varepsilon^2}\left(\phi^2 - 1\right)^2 + \frac{1}{2}\abs{\nabla \phi}^2.
\end{equation}

The corresponding evolution equation for the phase field $\phi$ reads
\begin{equation}
    \partial_t \phi = M \left(\Delta \phi - \frac{1}{\varepsilon^2}\left(\phi^3 - \phi\right)\right),
\end{equation}
where $M > 0$ is the mobility and $\varepsilon > 0$ controls the diffuse-interface width.

We studied this system on the unit square $\Omega = (0, 1)^2$ with periodic boundary conditions, i.e. again on a torus. For the circular benchmark used below, the initial data were chosen as
\begin{equation}
    \phi_0(x, y) =
    \begin{cases}
        1, & (x - 0.5)^2 + (y - 0.5)^2 < 0.2^2, \\
        -1, & \text{otherwise.}
    \end{cases}
\end{equation}

The parameters in the Firedrake script were set to $\varepsilon = 0.02$ and $M = 1$. Besides this circular seed, the script also offers a deterministic pseudo-random initialization in order to mimic the random-start educational examples.

Using a test function $v$, the weak form can be written as
\begin{equation}
    \left\langle \partial_t \phi, v \right\rangle + M\left\langle \nabla \phi, \nabla v \right\rangle + \frac{M}{\varepsilon^2}\left\langle \phi^3 - \phi, v \right\rangle = 0.
\end{equation}

In the implementation, the time derivative was approximated using the backward Euler method. This yields, at every time step, a nonlinear scalar problem
\begin{equation}
    \left\langle \frac{\phi_{i+1} - \phi_i}{\Delta t}, v \right\rangle + M\left\langle \nabla \phi_{i+1}, \nabla v \right\rangle + \frac{M}{\varepsilon^2}\left\langle \phi_{i+1}^3 - \phi_{i+1}, v \right\rangle = 0,
\end{equation}
which was solved in Firedrake using a \texttt{NonlinearVariationalProblem} and Newton's method.

Below we present a sequence of snapshots for the shrinking circular inclusion on the periodic square.

\begin{figure}[!h]
    \centering
    \begin{subfigure}[b]{0.2\textwidth}
        \includegraphics[width=\textwidth]{allen_cahn-circle-t0.pdf}
        \caption*{$t = 0.00$}
    \end{subfigure}
    \hfill
    \begin{subfigure}[b]{0.2\textwidth}
        \includegraphics[width=\textwidth]{allen_cahn-circle-t0.02.pdf}
        \caption*{$t = 0.02$}
    \end{subfigure}
    \hfill
    \begin{subfigure}[b]{0.2\textwidth}
        \includegraphics[width=\textwidth]{allen_cahn-circle-t0.04.pdf}
        \caption*{$t = 0.04$}
    \end{subfigure}

    \vskip\baselineskip

    \begin{subfigure}[b]{0.2\textwidth}
        \includegraphics[width=\textwidth]{allen_cahn-circle-t0.06.pdf}
        \caption*{$t = 0.06$}
    \end{subfigure}
    \hfill
    \begin{subfigure}[b]{0.2\textwidth}
        \includegraphics[width=\textwidth]{allen_cahn-circle-t0.08.pdf}
        \caption*{$t = 0.08$}
    \end{subfigure}
    \hfill
    \begin{subfigure}[b]{0.2\textwidth}
        \includegraphics[width=\textwidth]{allen_cahn-circle-t0.1.pdf}
        \caption*{$t = 0.10$}
    \end{subfigure}

    \caption{Snapshots from the Firedrake Allen-Cahn simulation on a unit square with periodic boundary conditions.}
    \label{fig:allen-cahn-evolution}
\end{figure}

\printbibliography

\end{document}
